\section{Planar and cylindrical setup}\label{c01-c1-planar-and-cylindrical-setup}

We begin with the positive path bounds that justify all later limits of finite cylinders. The same flux identity also verifies the critical normalization of the vertex weight.

Use the honeycomb lattice dual to the equilateral triangle tiling of side 1. Paths can begin/end at mid-edges ("ports"). A triangle visited has weight
\[
 x_*=(2\cos\lambda)^{-1},\qquad \lambda=\pi/8.
\]
Otherwise its weight is 1. Edges to/from terminal midpoints are traversed halfway; the paths counted are simple (in particular they cannot reverse immediately along an edge, nor repeat a lattice vertex (triangle)). Absolute SAW weights in what follows are products, without dividing out any partition function unless specified. A strip of height \(h\) (in layers, physical height \(h\sqrt3/2\)) has bottom and top on parallel tiling lines. Write \(A_h,b_h\) for the total weights from a given bottom port to bottom and to top respectively, inside the strip; the trivial path is excluded. All translated and reflected versions have the same totals.

On cuts formed by tiling lines we stop at the port just before leaving the allowed triangles. Centers of triangles are never on these cuts. The half-edge from a center up to (but not including) a port is inside that triangle, and the port is on its side but not at a tiling vertex. Thus a tiling line cutting a walk meets it through a perpendicular edge, not at a turn, and different nonparallel cuts never share a port.

The finite-domain identity below controls absolute path weights because convexity keeps all its phases in a fixed open half-plane.

\begin{lemma}[Critical flux]\label{lem:pv:flux}
For an outer boundary port of a finite convex domain bounded by tiling lines, start a path inwards and let it exit anywhere. Then
\[
 \sum_\gamma x_*^{|\gamma|} e^{i\sigma W(\gamma)}=1,\qquad \sigma=3/8
\]
where the sum is over the nontrivial walks to exit, \(|\gamma|\) counts vertices, and \(W\) is total turning angle along honeycomb edges. Every such walk has \(|W|\le\pi\), so its real phased weight is at least \(\cos(\sigma\pi)x_*^{|\gamma|}\).
\end{lemma}

\begin{proof}
At each interior lattice vertex take the sum of the phased weights of paths to its adjacent ports signed + when pointing away from that vertex and - when pointing towards it there (the initial empty path at the source points towards the interior). Sum also over vertices; internal terms cancel. Locally, paths pointing in towards an unused vertex each have two extensions out, exactly balancing them (\(2x_*\cos(\sigma\pi/3)=1\)).

The remaining negative terms occur when the vertex was used earlier; the tail describes one of the two ways round a polygon coming back towards that vertex on its unused edge. Pair them by reversing the traversal after the stem, thus changing the terminal port but keeping the same number of vertices. The stem lies outside the polygon, since it connects to the source without crossing it. Relative to the end of the stem just before the vertex, the two extra turning angles are \(\pm(2\pi-2\pi/3)\): on a positively oriented traversal the turn at the vertex would be \(+\pi/3\) (edge of the stem outside), and is replaced by \(-\pi/3\). Their phases therefore cancel. This uses only trivalence and the polygon turning-angle theorem.

Paths to the outer boundary are simple arcs in a convex polygon, so \(|W|\le\pi\): close them by the counterclockwise boundary arc and subtract from \(2\pi\) its intervening corner turns and the two right angles. In particular all real parts compare to absolute weights by the constant \(\cos(\sigma\pi)>0\).
\end{proof}

Arches to the right from a horizontal bottom side have \(W=-\pi\), those to the left \(W=\pi\), and opposite-side exits have \(W=0\). Passing from finite convex domains to an infinite strip gives the following normalization and preliminary bounds.

\begin{corollary}[Strip flux and crude bridge bounds]\label{cor:pv:strip-flux}
The strip masses \(A_h,b_h\) are finite and satisfy
\begin{equation}\label{eq:pv:strip-flux}
 b_h+\cos(\sigma\pi)A_h=1.
\end{equation}
Moreover,
\[
 A_{h+1}-A_h\le x_*^{-1} b_{h+1}^2,
 \qquad \frac{c}{h+1}\le b_h\le1.
\]
\end{corollary}

\begin{proof}
Cut the strip off at great lateral distances by tiling lines whose outward normals are rotated by \(\pm\pi/3\) from vertical upward. Lemma~\ref{lem:pv:flux} first gives finiteness of \(A_h,b_h\). To obtain equality in \eqref{eq:pv:strip-flux}, continue each side exit from its port into the part cut away, strictly increasing both height and outward projection, by a zigzag of vertical up and upward-outward edges to the top in the untruncated strip. The boundary edge crossing the cut is itself upward-outward. This costs a bounded factor at fixed \(h\), without repetition or unbounded multiplicity, since the first hit of the cut recovers the original exit. The weight of bridges so obtained tends to zero by finiteness, so the side-exit mass vanishes.

An arch visiting the last layer uses a vertex immediately below a top port: it enters the layer by the unique downward edge of a lower triangle there, and next goes to an upper triangle. Split there into two bridges sharing that vertex and ending at that port, and drop their avoidance. The first bridge starts at the fixed source; for each endpoint of the first, the reverse of the second can go to any bottom port, again with total weight bounded by \(b_{h+1}\). The shared vertex accounts for the factor \(x_*^{-1}\), proving the difference bound. By \eqref{eq:pv:strip-flux},
\[
 \frac1{b_{h+1}}-\frac1{b_h}
 \le x_*^{-1}\cos(\sigma\pi)\frac{b_{h+1}}{b_h}
 \le x_*^{-1}\cos(\sigma\pi).
\]
Summation gives the lower bound; the upper bound follows directly from the strip identity.
\end{proof}

\paragraph{Criticality of the weight.}
The unfolding argument follows the bridge-decomposition method of Hammersley and Welsh~\cite{HammersleyWelsh1962}, as used on the honeycomb lattice in~\cite[Section~3]{DuminilCopinSmirnov2012}. We include the geometric normalization needed for our ports. A vertex-to-vertex weak bridge of positive span (intermediate heights between those of its endpoints inclusively) can be extended at start and end to bracketing layer cuts using only strictly more extreme vertices, at most one extra vertex per end. Hence by the strip upper bound the weighted sum for each span and fixed start is uniformly bounded.

Count vertex SAWs of length \(L\) from a fixed start by splitting at a lowest vertex and considering the two weak half-space walks from there. On each, cut at the successive last maximum, last minimum of the remainder, and so on. This gives weak bridges of alternating sign with decreasing positive spans (the first two can be equal), in units of one third of a layer, with sum at most \(2L\) per half; there are no zero-height steps. There are \(O(\sqrt L)\) pieces and \(\exp(O(\sqrt L\log(L+1)))\) span lists. One can enumerate the two walks relative to the splitting vertex (two translation types suffice) and then the translate back to the given start is fixed. Thus multiply the uniform bounds allowing also a bounded weight factor per joint. SAW counts have exponential rate at most \(1/x_*\).

A strictly smaller rate would make the vertex SAW susceptibility (weighted sum over all lengths at \(x_*\)) finite, hence by trimming port half-edges would make \(\sum b_h\) finite, contradicting the bridge lower bound. The rate exists, e.g. by submultiplicativity and lattice transitivity. Thus the weight \(x_*\) is critical.

A flat cylinder identifies horizontal translates by an integer \(N>0\). The following fixed-width estimate will justify the row-transfer limits. Its constants may depend on \(N\).

\begin{lemma}[Fixed-width cylinder convergence]\label{lem:pv:cylinder-decay}
Fix \(N\). Let \(b_h^{(N)}\) and \(A_h^{(N)}\) denote the cylindrical versions of the strip masses from a fixed bottom port, with terminal port summed on the top and bottom cut respectively. Then \(\sum_h b_h^{(N)}<\infty\), the limit \(A_\infty^{(N)}\) is finite, and there are \(C_N,c_N>0\) such that
\[
 b_h^{(N)}+A_\infty^{(N)}-A_h^{(N)}
 \le C_N e^{-c_N h}.
\]
The same bounds apply when a terminal port is prescribed instead of summed.
\end{lemma}

\begin{proof}
A cylindrical bridge lifts to a plane strip bridge. Cut strip bridges at each interior horizontal layer line crossed just once. The weights of positive-height indecomposable pieces summed over all heights in the plane, from a fixed bottom start, total at most 1 by concatenation and \(b_h\le1\). More explicitly, all sequences concatenate freely by translations. If the total restricted to some finite height bound \(H\) exceeded 1, their weight after \(k\) concatenations would grow exponentially in \(k\), but would be bounded by \(\sum_{h\le kH}b_h\le kH\).

The total mass of pieces whose projections are simple on the cylinder is strictly less than 1: already in a single layer one can take a bottom-to-top bridge by following the internal zigzag arbitrarily far horizontally, and some such pieces repeat a vertex after projection. Summing the resulting geometric series bounds the total cylindrical bridge mass over all heights, after dropping avoidance between distinct pieces. Same-boundary arc weights are bounded by lifting and the uniform bound on planar arches.

It remains to strengthen summability to exponential decay. The sequences \(b_h^{(N)}\) and \(A_{h+1}^{(N)}-A_h^{(N)}\) have rational generating functions. Reveal layer by layer, retaining the occupied ports and the connectivity of dangling ends and prescribed terminals; keep only diagrams of paths, rejecting a closed cycle as soon as it forms. The update has finitely many states and constant coefficients, including periodic identification in each layer. A complete arc below the current line can be retained in a completed state. This counts single arcs with fixed or free terminal along the final cut, and also confined arcs starting and ending on the first cut. Translation by a lattice period in the layer direction identifies successive updates.

These sequences are nonnegative and summable. Rationality and absolute summability exclude poles on or inside the unit circle, so both sequences decay exponentially. Summing the second sequence over heights at least \(h\) gives the asserted same-boundary tail. Prescribing a terminal only reduces these positive sums.
\end{proof}
